\vspace*{4cm}
\begin{center}
    {\Huge\bfseries Résumé}
\end{center}
\vskip 1cm

\phantomsection
\addcontentsline{toc}{chapter}{Résumé}
\selectlanguage{french}

La transformation numérique du secteur de la santé a favorisé l'apparition de dispositifs médicaux connectés, de dossiers médicaux électroniques, de plateformes de télésurveillance et de systèmes d'aide à la décision fondés sur l'intelligence artificielle. Malgré la quantité importante de données générées, plusieurs systèmes eHealth restent fragmentés, statiques et faiblement personnalisés. La technologie du jumeau numérique offre une approche structurée pour répondre à ces limites en reliant des entités physiques, des processus ou des patients à des représentations virtuelles pouvant être mises à jour, analysées, simulées et utilisées pour soutenir la décision.

Ce mémoire de Master présente une revue structurée de l'état de l'art des jumeaux numériques appliqués à l'eHealth. Il étudie les fondements de l'eHealth et analyse l'évolution des définitions du jumeau numérique, depuis le modèle classique à trois dimensions reliant l'entité physique, l'entité virtuelle et la connexion, jusqu'aux approches étendues centrées sur les données et les services. Il étudie ensuite les niveaux de maturité et les applications dans le domaine de la santé, puis développe deux états de l'art complémentaires. Le premier concerne les jumeaux numériques en eHealth et en santé, notamment les jumeaux centrés sur le patient, les organes, les dispositifs médicaux, les processus hospitaliers, la santé publique, les services basés sur l'intelligence artificielle, l'interopérabilité, la confidentialité, l'éthique, la gouvernance et la validation. Le second se concentre sur les jumeaux numériques cardiovasculaires comme étude clinique représentative, en abordant l'électrophysiologie, l'hémodynamique, l'insuffisance cardiaque, la calibration patient-spécifique, les approches hybrides entre IA et modèles mécanistes, les jumeaux cardiaques à grande échelle et les défis de validation.

Le mémoire comprend une comparaison des définitions représentatives du jumeau numérique, deux tableaux comparatifs en orientation paysage et plusieurs figures de synthèse académique présentant les courants de définition, les rôles de l'IA, les catégories de jumeaux numériques en eHealth, les catégories cardiovasculaires, la maturité des preuves et la transition vers le PFE CardioTwin. Le travail déduit une architecture conceptuelle multicouche ainsi que des implications de conception pour CardioTwin, tout en conservant la portée générale du mémoire. CardioTwin est présenté comme une continuation pratique cardiovasculaire nécessitant une implémentation, une évaluation et une validation clinique futures.

\textbf{Mots-clés :} Jumeau numérique, eHealth, santé numérique, santé cardiovasculaire, intelligence artificielle, simulation, interopérabilité, aide à la décision clinique, CardioTwin.

\selectlanguage{english}
\clearpage
